\chapter*{Colour-Coded Thesis Review}
\addcontentsline{toc}{chapter}{Colour-Coded Thesis Review}
\markboth{Colour-Coded Thesis Review}{Colour-Coded Thesis Review}

This copy is an editorial and evidential review of the thesis. It carries two
kinds of mark. The first four colours appear as boxes and assess the section or
subsection that follows the corresponding heading. They do not alter the thesis
text, equations, tables, figures, or references.

\ReviewLegendItem{green}{GREEN}{The argument and presentation are sufficient
for the present thesis. Minor polishing may still be possible.}
\ReviewLegendItem{yellow}{YELLOW}{The content is relevant, but the explanation
is dense, repetitive, weakly supported, or harder to follow than necessary.}
\ReviewLegendItem{red}{RED}{The passage retains a generic, templated, or
textbook-restatement pattern that can read as machine-produced. This is a
stylistic assessment, not a claim about authorship or an AI-detector result.}
\ReviewLegendItem{grey}{GREY}{A reasoning step, evidence link, validation, or
required measurement is absent, uncertain, or inconsistent elsewhere.}

Green means sufficient, not perfect. A grey assessment does not make the
reported data invalid; it identifies where the current evidence does not fully
support the surrounding interpretation. Red is reserved for prose structure
rather than technical correctness. Where a section contains both sound work
and one important weakness, the more serious colour is used and the box names
the exact concern.

The remaining two colours record the state of the text rather than an
assessment of it, and are therefore applied to the words themselves.

\ReviewLegendItem{purple}{PURPLE}{The passage has been read, checked and
settled. Its wording is final and is not edited again.}
\ReviewLegendItem{blue}{BLUE}{Either new thesis text in a passage not yet
settled, or a comment raised against a settled passage.}

Purple is not a judgement that the passage is strong. It records that the
passage has been controlled and found correct, and that it already satisfies
what the thesis requires of it. Nothing inside it is reworded, tightened,
extended, or restructured afterwards. Rewording, sharpening, and stylistic
improvement are not grounds for changing a settled passage, and a revised
section that carries no blue at all is the normal case rather than an
oversight.

Where a settled passage appears to need something, the text is left as it
stands and the concern is raised beside it instead, as

\medskip
\Revised{The clearance gate holds the pose the tool reached and freezes the
descend clock, until a bare Enter releases it.}
\Comment{the frozen clock also governs the grind gate, which may be worth
stating here}
\medskip

\noindent A comment is a remark about the text, not a part of it. It appears
only in this annotated copy and never in the submitted thesis, so a settled
passage reads identically in both. It is raised only where the context or the
reasoning is genuinely incomplete without the addition, not for anything that
would merely read better.

Blue is used directly on the words themselves only in passages that have not
yet been settled, where it marks new text as a whole sentence or a whole word
and ends in a raised \textcolor{reviewbluemark}{$+$}. The two tints differ in
hue so that the boundary survives greyscale printing.

The clean submission files remain \texttt{Thesis.pdf} and
\texttt{Professor\_Draft.pdf}. This annotated copy is generated separately as
\texttt{Review\_Draft.pdf}.
